%-------------------------------------------------------------------------------
%                            BAB II  (VERSI MAGISTER KECERDASAN BUATAN / MKA)
%                       TINJAUAN KEPUSTAKAAN
%-------------------------------------------------------------------------------
% CATATAN: Adaptasi bab2.tex untuk program Magister Kecerdasan Buatan.
% Berkas asli (bab2.tex) TIDAK diubah. Pergeseran porsi:
%   - Geokimia + LIBS + fisika plasma DIKOMPRES menjadi satu seksi ringkas
%     "Domain dan Forward Model" (cukup untuk menjustifikasi inductive bias).
%   - Sisi AI DIPERLUAS: spectral source separation/unmixing, debat
%     interpretabilitas attention, taksonomi domain adaptation (sim-to-real),
%     physics-informed ML (level-arsitektur vs level-kerugian), efficient
%     attention, dan set/permutation-invariant encoding.
% Label persamaan diberi prefiks "bab2ai_" agar tidak bentrok dengan bab2.tex.
%-------------------------------------------------------------------------------

\chapter{TINJAUAN KEPUSTAKAAN}

% ======================================================================
% 2.1 Domain dan Forward Model (KOMPRES dari §2.1-§2.3 versi Fisika)
% ======================================================================
\section{Domain Studi Kasus dan Model Pembangkit Data}

Bagian ini meringkas latar fisis yang \emph{minimal namun cukup} untuk menjustifikasi \textit{inductive bias} arsitektural yang diusulkan; pembahasan geokimia dan instrumentasi yang lebih rinci berada di luar fokus tesis kecerdasan buatan ini.

\subsection{Karakteristik Spektrum LIBS sebagai Sinyal Campuran}
\textit{Laser-Induced Breakdown Spectroscopy} (LIBS) menghasilkan plasma mikro dari interaksi pulsa laser dengan permukaan material; saat plasma mendingin, atom dan ion tereksitasi memancarkan foton pada panjang gelombang diskret yang khas bagi tiap unsur \citep{legnaioli_laser-induced_2025, sawyers_database_2025}. Tanah vulkanik dipilih sebagai studi kasus karena komposisi multi-elemennya (Si, Al, Fe, Ca, Mg, Na, K) menghasilkan \textit{spectral overlap} parah dan efek matriks yang kuat---unsur transisi seperti Fe menyumbang ribuan garis emisi rapat---sehingga merepresentasikan skenario pemisahan sumber yang menantang sekaligus terkontrol \citep{harmon_laser-induced_2019, fabre_advances_2020, manelski_libs_2024}. Kompleksitas inilah, bukan aspek geokimianya, yang relevan bagi perancangan arsitektur.

\subsection{Forward Model Saha--Boltzmann sebagai Sumber Data Berlabel}

Pada kondisi \textit{Local Thermodynamic Equilibrium} (LTE) \citep{cristoforetti_local_2010}, populasi tingkat energi mengikuti distribusi Boltzmann
\begin{equation} \label{eq:bab2ai_boltzmann}
n_i = \frac{n_z \, g_i \exp(-E_i / k_B T_e)}{Z_z(T_e)},
\end{equation}
sedangkan keseimbangan antar tingkat ionisasi diatur Persamaan Saha
\begin{equation} \label{eq:bab2ai_saha}
\frac{n_e \, n_{z+1}}{n_z} = 2\,\frac{Z_{z+1}(T_e)}{Z_z(T_e)} \left(\frac{2\pi m_e k_B T_e}{h^2}\right)^{3/2} \exp\!\left(-\frac{E_{\mathrm{ion},z}}{k_B T_e}\right),
\end{equation}
dengan parameter atomik ($\lambda_{ki}, A_{ki}, E_k, g_k$) bersumber dari \textit{NIST Atomic Spectra Database} \citep{sawyers_database_2025}. Intensitas emisi tiap garis dihitung melalui koefisien emisi
\begin{equation} \label{eq:bab2ai_emiss}
j_\lambda = \frac{hc}{4\pi\lambda}\,A_{ki}\,n_k\,\phi_V(\lambda),
\end{equation}
dengan $\phi_V(\lambda)$ profil Voigt, yaitu konvolusi pelebaran Doppler (Gaussian)
\begin{equation} \label{eq:bab2ai_doppler}
\Delta\lambda_D = \lambda_{ki}\sqrt{\frac{8\ln 2\,k_B T_i}{m_k c^2}}
\end{equation}
dan pelebaran Stark (Lorentzian)
\begin{equation} \label{eq:bab2ai_stark}
\Delta\lambda_S \approx 2\,\omega_{ki}(T_e)\,\frac{n_e}{n_{e,\mathrm{ref}}} .
\end{equation}

Dari sudut pandang kecerdasan buatan, \textit{forward model} ini berperan sebagai \emph{generator data berlabel sempurna} yang tak terbatas dan terkendali. Sifat tersebut memungkinkan \textit{pre-training} berskala besar, \textit{ablation} terkontrol, dan evaluasi \textit{robustness} yang tak mungkin dilakukan hanya dengan data lapangan terbatas \citep{favre_towards_2025}.

\subsection{Struktur Generatif: Dekomposisi Linier Sumber--Campuran}
Inti yang menjadikan LIBS \textit{testbed} bernilai bagi penelitian ini adalah struktur generatifnya yang eksplisit: spektrum campuran (campuran) $S_{\mathrm{poly}}(\lambda)$ merupakan superposisi linier terbobot konsentrasi $c_z$ dari spektrum sumber (unsur tunggal) $S_{\mathrm{mono}}^{(z)}(\lambda)$ tiap unsur $z$:
\begin{equation} \label{eq:bab2ai_poly}
S_{\mathrm{poly}}(\lambda) = \sum_{z} c_z \, S_{\mathrm{mono}}^{(z)}(\lambda).
\end{equation}
Persamaan~\ref{eq:bab2ai_poly} adalah \textit{linear mixing model} (LMM) yang identik bentuknya dengan model yang melandasi \textit{spectral unmixing} pada citra hiperspektral \citep{palsson_hyperspectral_2018, ghosh_hyperspectral_2022}. Inilah jembatan konseptual yang memungkinkan persoalan kuantifikasi LIBS diformulasikan ulang sebagai persoalan pemisahan sumber pada ranah kecerdasan buatan, dan menjadi \textit{prior} struktural yang ditanamkan ke arsitektur (§2.4).

% ======================================================================
% 2.2 Pemisahan Sumber Spektral (BARU, sisi AI)
% ======================================================================
\section{Pemisahan Sumber dan \textit{Spectral Unmixing} Berbasis Pembelajaran}

Persoalan menguraikan sinyal campuran menjadi sumber-sumber penyusunnya beserta proporsinya merupakan problem klasik pembelajaran mesin. Pada \textit{spectral unmixing}, pendekatan berbasis \textit{autoencoder} mendominasi: struktur \textit{linear mixing model} secara implisit memberikan kendala arsitektural pada jaringan, sehingga \textit{decoder} berperan sebagai pustaka \textit{endmember} dan aktivasi laten sebagai \textit{abundance} \citep{palsson_hyperspectral_2018}. Ghosh et al.\ \citep{ghosh_hyperspectral_2022} memperluas paradigma ini dengan menggabungkan \textit{convolutional autoencoder} dan Transformer, memanfaatkan \textit{attention} untuk menangkap dependensi nonlokal antar pita spektral---arsitektur yang secara struktural paralel dengan yang diusulkan pada penelitian ini.

Perbedaan kunci penelitian ini terhadap literatur \textit{unmixing} klasik bersifat metodologis: sebagian besar metode \textit{unmixing} bersifat \emph{blind} (tanpa label \textit{abundance}) dan tidak memiliki \textit{ground truth} untuk memverifikasi sumber yang dipulihkan. Sebaliknya, kerangka LIBS dengan \textit{forward model} (Persamaan~\ref{eq:bab2ai_poly}) menyediakan label konsentrasi sempurna sekaligus pustaka sumber unsur tunggal yang diketahui, memungkinkan \textit{supervised source separation} dengan verifikasi eksternal yang objektif---kondisi yang langka pada domain \textit{unmixing} lain.

% ======================================================================
% 2.3 Kuantifikasi Klasik (RINGKAS, sebagai baseline)
% ======================================================================
\section{Metode Kuantifikasi Klasik sebagai \textit{Baseline}}

Sebagai pembanding, metode kuantifikasi LIBS konvensional dirangkum singkat. Kalibrasi univariat ($I = a c + b$) rentan terhadap efek matriks, sehingga pendekatan multivariat seperti \textit{Partial Least Squares Regression} (PLSR) mengeksploitasi seluruh spektrum untuk ketahanan yang lebih baik \citep{safi_multivariate_2018, zhang_partial_2020, wold_pls_2001}. \textit{Calibration-Free} LIBS (CF-LIBS) menurunkan konsentrasi langsung dari fisika plasma melalui \textit{Boltzmann plot} dan pelebaran Stark, namun akurasinya dibatasi validitas LTE, \textit{self-absorption}, dan ketidakpastian estimasi $T_e$/$n_e$ \citep{cristoforetti_local_2010, yang_self-absorption_2020, zeng_accurate_2023}. Favre et al.\ \citep{favre_towards_2025} menggabungkan \textit{forward model} dengan CNN untuk kuantifikasi \textit{real-time}, namun tetap meregresikan spektrum komposit secara langsung. Data referensi kebenaran (\textit{ground truth}) konsentrasi pada penelitian ini diperoleh dari \textit{X-ray Fluorescence} (XRF), metode standar analisis geokimia non-destruktif \citep{yamasaki_major_2023, ravansari_portable_2019}. Keempat pendekatan (PLS, CNN-\textit{only}, Transformer-\textit{only}, Informer) digunakan sebagai \textit{baseline} eksternal pada Bab~III.

% ======================================================================
% 2.4 Arsitektur Deep Learning (DIPERLUAS, inti AI)
% ======================================================================
\section{Arsitektur \textit{Deep Learning} untuk Spektrum Berdimensi Tinggi}

\subsection{\textit{Convolutional Neural Network} (CNN) 1D}
CNN satu dimensi efektif mengekstraksi fitur lokal dari spektrum---posisi puncak, lebar garis, struktur multiplet---melalui konvolusi \textit{kernel} geser pada sumbu panjang gelombang, dengan aktivasi ReLU dan \textit{Batch Normalisation} untuk stabilitas \citep{he_deep_2016, ioffe_batch_2015}. Dalam arsitektur yang diusulkan, CNN berperan sebagai \textit{tokenizer} spektral yang memproyeksikan spektrum mentah ke ruang laten sebelum pemodelan dependensi global.

\subsection{Transformer dan Mekanisme \textit{Self-Attention}}
Transformer \citep{vaswani_attention_2017} memodelkan relasi antar posisi spektral melalui
\begin{equation} \label{eq:bab2ai_attention}
\mathrm{Attention}(\mathbf{Q},\mathbf{K},\mathbf{V}) = \mathrm{softmax}\!\left(\frac{\mathbf{Q}\mathbf{K}^\top}{\sqrt{d_k}}\right)\mathbf{V}.
\end{equation}
\textit{Multi-Head Attention} memungkinkan model menangkap korelasi antar panjang gelombang dari berbagai subruang representasi secara paralel, sehingga mampu merekonstruksi konektivitas antar garis emisi satu unsur (mis.\ relasi multiplet dari tingkat ionisasi yang sama).

\subsection{Efisiensi \textit{Attention} pada Urutan Panjang}
Kompleksitas \textit{self-attention} konvensional yang $O(n^2)$ menjadi hambatan pada spektrum LIBS dengan $\sim$35.000 kanal. \textit{Informer} \citep{zhou_informer_2021} mereduksi kompleksitas menjadi $O(n \ln n)$ dengan mengeksploitasi \textit{sparsity} distribusi bobot \textit{attention} pada spektrum nyata, dan efektivitasnya pada spektrum geologi telah divalidasi \citep{walidain_informer-based_2026}. Pemilihan mekanisme \textit{attention} yang efisien karenanya bukan asumsi, melainkan keputusan desain yang harus dibandingkan secara eksplisit terhadap alternatif (Bab~III menjadikan Informer salah satu \textit{baseline}).

\subsection{Arsitektur \textit{Encoder--Decoder} dan \textit{Cross-Attention} sebagai \textit{Inductive Bias}}
Arsitektur \textit{Encoder--Decoder} memisahkan representasi masukan dari dekoding keluaran secara modular. Kontribusi metodologis penelitian ini adalah memetakan struktur generatif Persamaan~\ref{eq:bab2ai_poly} langsung ke topologi jaringan: \textit{encoder} memproses \emph{himpunan} spektrum sumber unsur tunggal $S_{\mathrm{mono}}^{(z)}$ menjadi memori spesifik-unsur, dan \textit{decoder} memproses spektrum campuran. Koneksi keduanya melalui \textit{Cross-Attention}, di mana \textit{query} $\mathbf{Q}_{\mathrm{dec}}$ dari spektrum campuran dicocokkan terhadap $\mathbf{K}_{\mathrm{enc}}, \mathbf{V}_{\mathrm{enc}}$ dari sumber, sehingga bobot \textit{cross-attention} dapat ditafsirkan sebagai \emph{koefisien pencampuran} $c_z$---konsisten dengan dekomposisi linier. Dengan demikian fisika ditanamkan pada \emph{level arsitektur}, berbeda dari paradigma \textit{physics-informed neural networks} yang umumnya menanam fisika pada \emph{suku fungsi kerugian} \citep{karniadakis_physics-informed_2021, raissi_physics-informed_2019}. Karena \textit{encoder} memproses himpunan sumber yang \emph{tak berurutan}, desain yang invarian terhadap permutasi relevan, sebagaimana diformalkan pada kerangka \textit{Deep Sets} \citep{zaheer_deep_2017}.

\subsection{Interpretabilitas \textit{Attention} dan Debat Kesetiaan (\textit{Faithfulness})}
Penafsiran bobot \textit{attention} sebagai penjelasan model masih diperdebatkan. Jain dan Wallace \citep{jain_attention_2019} berargumen bobot \textit{attention} tidak konsisten dengan ukuran kepentingan fitur lain dan dapat diganti tanpa mengubah prediksi, sehingga "bukan penjelasan". Wiegreffe dan Pinter \citep{wiegreffe_attention_2019} membantah bahwa kesimpulan tersebut bergantung pada definisi penjelasan dan mengusulkan uji diagnostik yang lebih ketat. Bibal et al.\ \citep{bibal_attention_2022} merangkum bahwa \textit{attention} dapat menjadi penjelasan pada kondisi tertentu, terutama bila tersedia acuan eksternal untuk memverifikasinya. Keterbatasan utama seluruh debat ini adalah \emph{ketiadaan \textit{ground truth} objektif} tentang apa yang "seharusnya" di-\textit{attend}. Penelitian ini memanfaatkan keunikan domain LIBS: posisi garis emisi NIST menyediakan \textit{oracle} fisik objektif, sehingga kesetiaan bobot \textit{cross-attention} dapat diukur secara kuantitatif---sebuah kontribusi langsung pada perdebatan ini.

\subsection{Adaptasi Domain dan Transfer \textit{Sim-to-Real}}
Model yang dilatih pada data sintetis (\textit{forward model} LTE) mengalami penurunan kinerja pada data eksperimen akibat \textit{domain shift}: efek instrumentasi, \textit{continuum background} residual, dan deviasi dari LTE ideal. Adaptasi domain adversarial menjawab hal ini; \textit{Domain-Adversarial Training} via \textit{Gradient Reversal Layer} (GRL) \citep{ganin_domain_2016} membalik gradien klasifier domain agar \textit{encoder} menghasilkan fitur \textit{domain-invariant}. Alternatif lain mencakup \textit{Adversarial Discriminative Domain Adaptation} (ADDA) yang memisahkan \textit{encoder} sumber dan target \citep{tzeng_adversarial_2017}. Pemilihan GRL pada penelitian ini didasarkan pada kompatibilitasnya dengan pelatihan \textit{end-to-end} satu tahap; perbandingan eksplisit terhadap kondisi tanpa GRL menjadi bagian \textit{ablation} (Bab~III). Kerangka \textit{sim-to-real}---\textit{pre-training} pada simulasi fisika lalu \textit{fine-tuning} pada data nyata terbatas---merupakan strategi \textit{low-resource learning} yang berdiri sendiri sebagai kontribusi.

% ======================================================================
% 2.5 Evaluasi dan Metrik (DIPERLUAS dengan faithfulness)
% ======================================================================
\section{Evaluasi Model dan Metrik Kesetiaan}

\subsection{Metrik Regresi dan Kualitas Spektrum}
Performa kuantifikasi dievaluasi dengan RMSE, $R^2$, dan MAPE per unsur. Kualitas spektrum eksperimen dinilai melalui \textit{Signal-to-Noise Ratio}
\begin{equation} \label{eq:bab2ai_snr}
    \mathrm{SNR} = \frac{I_{\text{peak}} - I_{\text{bg}}}{\sigma_{\text{bg}}},
\end{equation}
dengan ambang $\mathrm{SNR} \geq 10$ sebagai kriteria kelayakan \citep{khan_magnetic_2022, guerrini_optimization_2025}.

\subsection{Validasi Posisi Garis terhadap NIST sebagai \textit{Oracle}}
Deteksi puncak dan pencocokan garis terhadap \textit{NIST ASD} menyediakan acuan posisi $\lambda_{\mathrm{NIST}}$ yang objektif \citep{legnaioli_laser-induced_2025}. Selain memvalidasi \textit{forward model}, acuan ini berfungsi sebagai \textit{ground truth} untuk menilai apakah bobot \textit{cross-attention} memang terkonsentrasi pada garis emisi unsur yang benar---menjembatani evaluasi numerik dengan evaluasi kesetiaan interpretabilitas (lihat §2.4.5). Robustness terhadap perturbasi realistis (\textit{Gaussian/Poisson/spike noise}, \textit{wavelength shift}, \textit{instrumental broadening}, \textit{baseline drift}, \textit{missing-line suppression}) dievaluasi secara sistematis sebagai uji ketahanan model terhadap \textit{distribution shift}.
